\pdfoutput=1
\documentclass[11pt]{article}

\usepackage[margin=1in]{geometry}
\usepackage{amsmath,amssymb}
\usepackage{booktabs}
\usepackage{graphicx}
\usepackage{hyperref}
\usepackage{xcolor}
\usepackage{lineno}
\usepackage[numbers,sort&compress]{natbib}

\newcommand{\IIA}{\mathrm{IIA}}
\newcommand{\Gr}{\mathrm{Gr}}
\newcommand{\R}{\mathbb{R}}
\newcommand{\dgr}{d_{\Gr}}
\newcommand{\TODO}[1]{\textcolor{red}{[#1]}}
\renewcommand{\arraystretch}{1.15}

\title{Distributed Alignment Search Recovers a Subspace, Not the Subspace}

\author{Anonymous}
\date{}

\begin{document}
\linenumbers
\maketitle

\begin{abstract}
Applications of Distributed Alignment Search report \emph{a} subspace as though it
were \emph{the} subspace for a task. We show it is not, in a setting clean enough
to measure: grokked modular arithmetic transformers, where the causal variable is
known and recoverable. Ten independently trained models all reach $\IIA$ between
$0.94$ and $1.00$, and the subspaces recovered from them have pairwise overlap
$\approx 0.008$ against a chance value of $k/d = 0.016$ for random
$k$-dimensional subspaces of $\R^{128}$. The recovered subspaces are less aligned
than chance, and pairwise geodesic distances on $\Gr(2,128)$ occupy the narrow band
$[1.93, 2.19]$ radians, close to the maximum $\pi/\sqrt{2} \approx 2.22$. Every one
supports equally valid interventions, and rotating a solution by more than a
radian barely changes accuracy. What is invariant is not the subspace but the
geometry inside it: label centroids lie on a circle and equivariance under the
operation's symmetry group holds in every solution. We argue that alignment
results should report such invariants rather than basis vectors.
\TODO{The claim currently cannot separate model-level from fitting-level
degeneracy; see \S\ref{sec:decomp}. Resolve before submission.}
\end{abstract}

\section{Introduction}

\TODO{Frame: DAS papers report the subspace for IOI, for arithmetic, for board
games. If the object is only defined up to a large equivalence class, reporting
one representative and interpreting its basis vectors is not meaningful. We test
this where the ground truth is cleanest.}

\section{Background}

\TODO{Short. The Grassmannian $\Gr(k,d)$, principal angles, geodesic distance
$\dgr(Q_1,Q_2) = \sqrt{\sum \theta_i^2}$, DAS as optimisation over $\Gr(k,d)$.
This section exists only to make ``overlap $0.008$ against chance $0.016$''
interpretable, and it earns its place here in a way it would not elsewhere: a
point on $\Gr(k,d)$ is exactly ``a subspace up to rotation'', and the finding is
that grokked models do not select a point.}

\section{Results}

\subsection{Subspace degeneracy}

Ten independently trained grokked multiplication models, DAS fitted to each:

\TODO{Table: pairwise overlap $\approx 0.008$ vs chance $0.016$; geodesic
distances in $[1.93, 2.19]$ radians against a maximum of $2.22$; per-model $\IIA$
in $[0.94, 1.00]$. Plus the basin result: rotating a solution more than a radian
on $\Gr(k,d)$ barely reduces $\IIA$.}

The sub-chance alignment is the part that needs an account. Independently fitted
subspaces being unrelated would give overlap at chance; being \emph{below} chance
suggests the solutions are actively spread rather than merely unconstrained.
\TODO{We do not have an explanation. State that plainly rather than speculating.}

\subsection{Whether the degeneracy is in the model or in the fitting}
\label{sec:decomp}

\TODO{GATING EXPERIMENT, NOT YET RUN.
The ten-seed result varied the model seed and fitted DAS once per model, so model
non-uniqueness and optimisation sensitivity are confounded. Three conditions
separate them:
(A1) one model, repeated fits, delta-PCA initialisation --- measures batch-sampling
sensitivity only, since the initialisation is a deterministic function of the data;
(A2) one model, repeated fits, random orthogonal initialisation --- measures
whether the model pins a subspace at all;
(B) ten models, one fit each --- the current claim.
If A2 gives near-orthogonal subspaces on a single model, the degeneracy is in the
optimisation landscape and the paper's claim changes accordingly.
Script: \texttt{experiments/degeneracy\_decomposition.py}.}

\subsection{The geometry is invariant even though the coordinates are not}

\TODO{The load-bearing counterpoint, and what keeps the paper from being purely
negative. In every recovered subspace: label centroids lie on a circle, and
intervention-based equivariance under the operation's symmetry group holds.
Coordinates are arbitrary; the structure inside them is not.}

\subsection{Candidate invariants}

\TODO{What to report instead of a basis: circle geometry of label centroids,
equivariant fraction, intrinsic dimension $k^*$, and the weight-space subspace.
Weight-space SVD reaches $\IIA = 0.999$ without optimisation against DAS's
$0.150$, but the comparison is not dimension-matched ($k=32$ against $k=8$) and
cannot be reported as written. If SVD subspaces agree across models where DAS
subspaces do not, the result becomes ``DAS is degenerate, weight geometry is
not'', which is considerably stronger.}

\subsection{Subspace dimension}

Corrected sweep, three seeds, random orthogonal initialisation matching the
reference implementation, with the warm-started variant reported separately:

\begin{table}[t]
\centering\small
\setlength{\tabcolsep}{4pt}
\begin{tabular}{lccccccc}
\toprule
& $k{=}1$ & $k{=}2$ & $k{=}4$ & $k{=}8$ & $k{=}16$ & $k{=}32$ & $k{=}64$ \\
\midrule
DAS (random init)    & 0.009 & 0.063 & 0.291 & 0.722 & 0.919 & 0.957 & 0.980 \\
DAS (delta-PCA init) & 0.017 & 0.074 & 0.365 & 0.806 & \textbf{1.000} & 1.000 & 1.000 \\
NL-DAS               & 0.185 & 0.517 & 0.711 & 0.909 & 0.852 & 0.844 & 0.922 \\
Structured VAE       & 0.957 & \textbf{1.000} & 1.000 & 1.000 & 1.000 & 1.000 & 1.000 \\
\bottomrule
\end{tabular}
\caption{Interchange accuracy on grokked modular addition, mean over three seeds.
\TODO{Add CI row.} Warm-starting DAS from the top-$k$ right singular vectors of
the intervention deltas outperforms random orthogonal initialisation at every $k$,
reaching $1.000$ at $k=16$ where the standard initialisation does not saturate
even at $k=64$. Initialisation is therefore not incidental to DAS.}
\end{table}

\TODO{Decide whether the structured VAE row belongs in this paper at all. The
split plan says no; it appears here only because it is in the same sweep.}

\subsection{Where the degeneracy appears}

\TODO{The atlas, used as evidence for this claim only rather than as a survey:
grokked models have the circular invariant with degenerate coordinates,
non-grokked models have neither. Grokking rate scales with data ($8\%$ at $p=53$,
$71\%$ at $p=211$) and depth is non-monotonic, with two layers unlocking
operations four layers lose. Memorisation artifacts belong here too: squaring and
cubing reach $\IIA = 0.857$ at $k=2$ without grokking.}

\subsection{Stochastic grokking}

\TODO{Composite addition groks at 5 of 10 seeds, power at 0 of 10. Same
``the seed decides'' theme as the degeneracy result, and independently solid.
Single-condition classifications near the boundary are unreliable.}

\subsection{Equivariance requires a group action, not linearity}

\TODO{Affine ($2a + 3b + 5 \bmod 113$) is linear and fails equivariance at $0\%$
across four primes: under $a \mapsto a+g$ the output shifts by $2g$. Cubic sum
passes where cubing alone does not; max is piecewise equivariant. Operation
linearity and causal-variable linearity are different properties, which is why
the group-structured operations are the ones DAS handles.}

\section{Discussion}

\TODO{Implications for alignment practice: report invariants, report seed spread,
do not interpret basis vectors. Note that at $k = d$ the projection is the
identity and any linear method attains $\IIA = 1$ trivially, so the question is
always how far below $d$ a method reaches.}

\paragraph{Limitations.}
\TODO{Synthetic modular arithmetic at $d_\text{model} = 128$. The
model-versus-fitting decomposition is not yet run. The weight-space comparison is
not dimension-matched. Atlas cells near the grokking boundary flip across seeds
and the atlas table is single-run.}

\section{Methods}

\TODO{Grokking training. DAS fitting and initialisation. Equivariance test.
Subspace overlap and geodesic distance. The chance baseline $k/d$ and the null
distribution of random subspaces on $\Gr(k,d)$.}

\section*{Data and Code Availability}
\TODO{Repository and archive.}

\bibliographystyle{plainnat}
\bibliography{references,new_bib_entries}
\end{document}
